\begin{table}[htbp]
\centering
\caption{The study plan, from the prediction rung that releases nothing to the joint inversion of the whole campaign across all three filters (Filter~1 AR6, Filter~2 BS45, Filter~3 BS17) and two reference single layers. Residuals are rms departures in percent of reflectance or of relative transmittance, one per measurement channel. Rung~0 is a direct prediction evaluated with reference-transferred optical constants, not a fit: it adjusts zero parameters and tests parameter-free predictive agreement. A rung is a kind of study rather than a step in a sequence, so one rung can carry more than one component: rung~3 is listed twice, once for each beam splitter. Rung~4, the beam-aperture study on Filter~3 (\texttt{06\_BS17\_aperture}), and the unpolarized variant of Filter~2 (\texttt{03\_BS45\_unpolarized}) are deposited beside the baseline ladder and discussed in the text rather than printed here.}
\resizebox{\linewidth}{!}{%
\begin{tabular}{lrrrp{0.20\linewidth}p{0.30\linewidth}}
\hline
Rung & Free & Points & Pts/par & Released blocks & rms residual per measurement (\%) \\
\hline
0. prediction & 0 & 2266 & --- & none & reference SiO2 0.10; reference Nb2O5 0.16; Filter 1 0.55; Filter 2 s 1.10; Filter 2 p 0.79; Filter 3 s 1.57; Filter 3 p 1.20 \\
1. the floor & 2 & 516 & 258 & 2 thicknesses & reference SiO2 0.10; reference Nb2O5 0.15 \\
2. one coating & 6 & 250 & 42 & 6 thicknesses & Filter 1 0.11 \\
3. polarization resolved & 16 & 500 & 31 & 16 thicknesses & Filter 2 s 0.37; Filter 2 p 0.25 \\
3. polarization resolved & 17 & 1000 & 59 & 17 thicknesses & Filter 3 s 0.96; Filter 3 p 0.48 \\
5. cross-checking & 41 & 2266 & 55 & 41 thicknesses & reference SiO2 0.10; reference Nb2O5 0.15; Filter 1 0.11; Filter 2 s 0.37; Filter 2 p 0.25; Filter 3 s 0.96; Filter 3 p 0.48 \\
\hline
\end{tabular}%
}
\label{tab:ladder}
\end{table}
